\chapter{Higher-Order Systems}
\label{ch:highorder}
Dynamic systems can have more than two energy stores or modes: a multistage feedback
loop, a multi-pole amplifier, or a feed line with infinitely many modes. For the
stable all-pole responses developed here, little new machinery is required: they
factor into the first- and second-order pieces of \cref{ch:linsys}, and their
frequency responses add in decibels and degrees. Chapter~\ref{ch:filters} later
uses the same structure to explain filter goals and named filter families.

This chapter derives \(-20n\)~dB/decade, identifies a three-lag feedback limit,
and shows how pure delay and infinitely many modes extend the finite-order picture.

\section{Order Is the Number of Independent Energy Stores}
The counting rule is the same one \cref{ch:linsys} used, just not stopped at two. A
system with \(n\) independent energy stores has \(n\) state variables, so its state
matrix \(A\) is \(n\times n\), its characteristic polynomial
\(\det(sI-A)\) has degree \(n\), and it therefore has exactly \(n\) poles:
\[
\dot{\state}=A\state+Bu,
\qquad
\state\in\mathbb{R}^n
\qquad\Longrightarrow\qquad
n \text{ poles} .
\]
Nothing about \(n=5\) is harder than \(n=2\) in principle. Factoring a fifth-degree
polynomial by hand is unpleasant, which is why the next section matters.

\section{Factoring Stable All-Pole Systems}
\label{sec:factoring}
Here is the structural fact that makes higher order tractable.

A physical model with real-valued parameters---resistances, inductances, and
masses, for example---has a characteristic polynomial with real coefficients. Its
complex roots therefore occur in \textbf{conjugate pairs}: if
\(\sigma+\jj\omega\) is a root, so is \(\sigma-\jj\omega\). (Otherwise multiplying the
factors out would leave imaginary coefficients.) For the stable all-pole responses
considered here, the denominator factors into
\begin{itemize}
\item \textbf{real linear factors} \((s+1/\tau_k)\), one per real pole, and
\item \textbf{irreducible real quadratic factors}
\(\bigl(s^2+2\zeta_k\omega_{0k}s+\omega_{0k}^2\bigr)\), one per conjugate pair.
\end{itemize}
Thus the stable all-pole transfer function considered in this section can be written
as a \emph{product}:
\[
\boxed{\;
G(s)=K\prod_k\underbrace{\frac{1}{1+\tau_ks}}_{\text{first order}}
\prod_j\underbrace{\frac{\omega_{0j}^2}{s^2+2\zeta_j\omega_{0j}s+\omega_{0j}^2}}_{\text{second order}} \;}
\]
An even-order system is \(n/2\) quadratic sections; an odd-order system is
\((n-1)/2\) quadratics plus one real pole. Numerator factors add zeros; their
filter-design role is introduced in \cref{ch:filters}.

\begin{controlsbox}
\textbf{Factorization makes the order readable.} A fifth-order all-pole response is
two second-order sections and one first-order section in series. Its poles are the
combined poles of those three sections.
\end{controlsbox}

\section{Cascades Add in Decibels and in Degrees}
\label{sec:cascadeadd}
Now the frequency-domain payoff. \Cref{ch:linsys} showed that cascaded blocks
\emph{multiply}, \(G=G_1G_2\). Multiplication of complex numbers multiplies magnitudes
and adds angles, so
\[
|G(\jj\omega)|=|G_1(\jj\omega)|\cdot|G_2(\jj\omega)|,
\qquad
\angle G=\angle G_1+\angle G_2 .
\]
Take \(20\log_{10}\) of the magnitude equation and the product becomes a sum, because
that is what logarithms do:
\[
\boxed{\;
20\log_{10}|G|=20\log_{10}|G_1|+20\log_{10}|G_2|,
\qquad
\angle G=\angle G_1+\angle G_2 \;}
\]
\textbf{This is the reason Bode plots use decibels.} On a dB scale, cascading is
addition. A high-order Bode plot is the sum of its sections' plots, and the sketching
rules of \cref{ch:splane} (each pole bends the slope down by
\SI{20}{\decibel}/decade, each zero up) are simply this addition performed
graphically.

\subsection{Deriving the \texorpdfstring{\(-20n\)}{-20n} dB/Decade Roll-Off}
The rule quoted throughout the book now follows by adding the one-pole results of
\cref{sec:asymptotes}. Well above its corner, each pole contributes
\(-\SI{20}{\decibel}\) per decade and approaches \(-90^\circ\) of phase lag. A
second-order section contains two poles, so it contributes twice those amounts. For
\(n\) poles total, the magnitude contributions add:
\[
\underbrace{(-20)+(-20)+\cdots}_{n\text{ terms}}
=-20n\ \si{\decibel}\text{ per decade}.
\]
The phase lag likewise approaches \(-90n^\circ\) far above the corners. Both facts
are consequences of \cref{sec:cascadeadd}.

\begin{workedbox}[: three identical lags]
Cascade three identical first-order stages, each with corner \(\omega_0\):
\[
G(s)=\frac{1}{(1+s/\omega_0)^3}.
\]
Far above \(\omega_0\), the magnitude falls at \(-\SI{60}{\decibel}\) per decade and
the phase approaches \(-270^\circ\). At the corner itself each stage is down
\SI{3.01}{\decibel} and lagging \(45^\circ\), so the cascade is down
\(3\times3.01\approx\SI{9}{\decibel}\) and lagging \(135^\circ\)---all three numbers
obtained by addition, with no new algebra. Keep this system in mind; the next section
shows it is the one that can oscillate.
\end{workedbox}

\section{Why Three Stable Lags Can Destabilize a Feedback Loop}
\label{sec:threepoles}
For the stable, strictly proper cascades of first-order lags in ordinary negative
feedback considered here, \cref{sec:cascadeadd} gives a useful phase-accounting
limit.

Recall the condition for closed-loop poles on the \(\jj\omega\) axis:
\(L(\jj\omega)=-1\), meaning \(|L|=1\) and \(\angle L=-180^\circ\) at the \emph{same}
frequency. Since each first-order lag contributes at most \(-90^\circ\):
\begin{description}
\item[One pole] the phase approaches \(-90^\circ\) but never reaches
\(-180^\circ\). The condition is unreachable at any gain in this model, as
\cref{ch:feedback} found by moving the pole directly.
\item[Two poles] the phase approaches \(-180^\circ\) only asymptotically, as
\(\omega\to\infty\), where the magnitude has already fallen to zero. So \(|L|=1\) and
\(\angle L=-180^\circ\) never coincide. This two-lag model remains stable, though
its poles can be driven arbitrarily close to the axis (which is why it can ring hard
without oscillating).
\item[Three poles] the phase reaches \(-180^\circ\) at a \emph{finite} frequency, where
the magnitude is still nonzero. With sufficient gain, both conditions can coincide
and a pole pair can reach the imaginary axis.
\end{description}

For the three-lag cascade of the box above, put \(L(s)=K/(1+s/\omega_0)^3\). The phase
is three times a single lag's:
\[
\angle L=-3\arctan\!\frac{\omega}{\omega_0}=-180^\circ
\quad\Longrightarrow\quad
\arctan\frac{\omega}{\omega_0}=60^\circ
\quad\Longrightarrow\quad
\omega_{\mathrm{osc}}=\sqrt3\,\omega_0 .
\]
At that frequency each stage has magnitude
\(1/|1+\jj\sqrt3|=1/2\), so the cascade has magnitude \(K/8\). Setting \(|L|=1\):
\[
\boxed{\;K=8\quad\text{puts the poles exactly on the }\jj\omega\text{ axis at }
\omega=\sqrt3\,\omega_0 . \;}
\]
Below \(K=8\) the loop is stable with positive margins; above it, the pole pair is in
the right half-plane. This single system is the one plotted three ways in
\cref{fig:fb-margins} (\(K=3\), healthy margins),
\cref{fig:opamp-rootlocus} (the poles migrating, crossing at \(K=8\)), and
\cref{fig:oscillator-loopgain} (\(K=8\), margins exactly zero)---and now the arithmetic
behind all three lives here.

\section{Pure Delay: Phase Lag Without Any Poles}
\label{sec:puredelay}
A delay of \(T\) seconds is not a pole at all, yet it destabilizes loops more
effectively than poles do. Its transfer function is
\[
G(s)=\ee^{-sT},
\qquad
|G(\jj\omega)|=1,
\qquad
\angle G(\jj\omega)=-\omega T .
\]
Interpret those two facts together: the magnitude is \emph{never} attenuated, while the
phase lag grows \emph{without bound} and linearly in \(\omega\). So a delay crosses
\(-180^\circ\), then \(-540^\circ\), then \(-900^\circ\), forever, always at full
amplitude. A real delay therefore supplies whatever additional phase a loop needs;
the loop oscillates only if its remaining gain can also satisfy \(|L|=1\) at one of
those phase crossings. This is why transport lag is the classic difficulty in process
control. A matched transmission line (\cref{ch:lines}) supplies propagation delay;
reflections at its boundaries create the standing-wave resonances.

\section{The Limit of Infinitely Many Poles}
\label{sec:infinitepoles}
Push the order up without bound and something qualitatively new appears. A cascade of
\(n\) identical lags, each of time constant \(T/n\), tends to a pure delay of \(T\) as
\(n\to\infty\)---infinitely many poles conspiring to produce exactly the
\(\ee^{-sT}\) of the delay section. That is not an idle limit: it is what a
transmission line \emph{is}. \Cref{ch:lines} models a line as an infinite ladder of
infinitesimal \(L\) and \(C\) sections, finds a wave equation rather than an ODE, and
sees the infinite family of resonances directly in the periodic poles of
\(\tan\beta\ell\). The lumped machinery of Part~II reaches its edge there, and the
transition is continuous.

\begin{exambox}
The exam does not ask for fifth-order pole coordinates. It asks for the roll-off of an
\(n\)-pole response (\(-20n\)~dB/decade, \cref{sec:cascadeadd}), why too much loop
gain can oscillate (\cref{sec:threepoles}), and why a transmission line requires a
distributed rather than lumped model. These are all consequences of one idea: a
high-order system is built from first- and second-order sections, added in decibels
and degrees, until the distributed limit takes over.
\end{exambox}
